\section*{Bab \ref{ch:b3:heterocycles} --- Kimia Heterosiklik}

\begin{solution}{exo:b3:heterocycles:1}
Oksazola: O menyumbang dua (satu pasangan elektron bebas dalam sistem $\pi$, satu di bidang), N satu, dan
ketiga karbon tiga: 6. Tiazola: sama dengan S: 6. Pirimidina: satu dari setiap
atom dari keenam atom, dengan kedua pasangan elektron bebas nitrogen di bidang: 6. Purina: sembilan atom,
dengan nitrogen N--H menyumbang dua dan yang lain masing-masing satu: 10, sebuah bilangan Hückel
($4n + 2$, $n = 2$).
\end{solution}

\begin{solution}{exo:b3:heterocycles:2}
Piperidina (11.1) $>$ DMAP (9.6) $>$ imidazola (7.0) $>$ piridina (5.2) $\gg$ pirola (sekitar $-4$,
terprotonasi pada karbon).
\end{solution}

\begin{solution}{exo:b3:heterocycles:3}
2-Bromopirola (pirola begitu reaktif sehingga mudah terpolibrominasi);
2-bromofuran; 3-bromoindola.
\end{solution}

\begin{solution}{exo:b3:heterocycles:4}
2,5-Dimetilfuran; 1,2,5-trimetilpirola; 2,5-dimetiltiofena.
\end{solution}

\begin{solution}{exo:b3:heterocycles:5}
2-Metoksipiridina. Adisi metoksida pada C2 menempatkan muatan negatif
\omterm{def:b3:heterocycles:snar}{kompleks Meisenheimer} pada nitrogen; pada C3, muatan itu hanya dapat tersebar pada karbon-karbon, sehingga
zat antara, dan keadaan transisi sebelumnya, jauh lebih tinggi energinya.
\end{solution}

\begin{solution}{exo:b3:heterocycles:6}
2-Aminopiridina (sebagai garam natriumnya, lalu terprotonasi pada pengerjaan). Ion amida
teradisi pada C2, sehingga terbentuk hasil adisi $\sigma$ anionik dengan H dan $\ce{NH2}$ pada C2 dan muatan negatif
pada nitrogen cincin; hasil adisi itu melepaskan hidrida, yang mengambil sebuah proton dari
gugus amino dan lepas sebagai $\ce{H2}$.
\end{solution}

\begin{solution}{exo:b3:heterocycles:7}
Oksigen N-oksida mengembalikan sebuah pasangan elektron bebas ke dalam cincin: struktur-struktur
resonansinya membawa muatan negatif pada C2, C4, dan C6. Elektrofil menyerang C4 (C2
dekat dengan nitrogen yang bermuatan positif). Oksigennya lalu disingkirkan dengan
fosfor triklorida, sehingga diperoleh 4-nitropiridina.
\end{solution}

\begin{solution}{exo:b3:heterocycles:8}
2-Hidroksipiridina (OH pada C2) dan 2-piridon (N--H dan C=O). Piridon mempertahankan enam
elektron $\pi$ dalam cincin: struktur resonansi zwitterioniknya, dengan $\mathrm{O^-}$ dan
$\mathrm{N^+}$ yang berbagi pasangan elektron bebasnya, adalah sebuah piridinium-olat; piridon adalah amida yang
pasangan elektron bebas nitrogennya melengkapi sekstet, dan ikatan hidrogen C=O dan N--H yang kuat
menjadikannya lebih disukai di dalam air.
\end{solution}

\begin{solution}{exo:b3:heterocycles:9}
2-Nitrobenzaldehida, dua ekuivalen metil asetoasetat, dan amonia: penghambat
kanal kalsium nifedipina.
\end{solution}

\begin{solution}{exo:b3:heterocycles:10}
Ena-hidrazina ke arah gugus $\ce{CH2}$ (dalam, lebih tersubstitusi) memberikan
2,3-dimetilindola; ena-hidrazina ke arah gugus metil (ujung) memberikan 2-etilindola.
Asam lemah terutama memberikan yang pertama, melalui ena-hidrazina yang lebih stabil;
asam yang lebih kuat menaikkan proporsi 2-etilindola.
\end{solution}

\begin{solution}{exo:b3:heterocycles:11}
Dalam $\mathrm S_\mathrm NAr$, tahap penentu lajunya adalah adisi, yang tidak memutus ikatan
C--X; fluor yang lebih elektronegatif lebih menstabilkan keadaan transisi anionik
dan \omterm{def:b3:heterocycles:snar}{kompleks Meisenheimer}. Dalam $\mathrm S_\mathrm N2$, ikatan C--X putus dalam
keadaan transisi, dan ikatan C--I yang lemah lepas paling cepat.
\end{solution}

\begin{solution}{exo:b3:heterocycles:12}
Piridina: tiga sinyal dengan perbandingan 2\,:\,1\,:\,2, yang pertama di dekat \qty{8.6}{ppm}. Pirola: dua sinyal
yang sama besar pada 6.7 dan \qty{6.2}{ppm} serta sinyal N--H yang lebar di dekat \qty{8.1}{ppm}. Furan: dua sinyal yang sama besar pada
7.4 dan \qty{6.4}{ppm}, tanpa N--H. Tiofena: dua sinyal yang sama besar pada 7.33 dan \qty{7.12}{ppm}, yang saling
berdekatan, di antara kedua sinyal furan.
\end{solution}

\begin{solution}{pb:b3:heterocycles:1}
\textbf{1.} $\ce{C6H5NHNH2 + (CH3)2CO -> C6H5NHN=C(CH3)2 + H2O}$.

\textbf{2.} Nitrogen amina teradisi pada karbon karbonil, lalu air lepas, seperti pada
pembentukan imina; asam mengaktifkan gugus karbonil dan mengubah OH hasil adisinya
menjadi gugus pergi.

\textbf{3.} $\ce{NH2}$ ujung: pasangan elektron bebas nitrogen yang lain terkonjugasi dengan
cincin fenil dan lebih terhalang.

\textbf{4.} $\qty{108.0}{g/mol}$; $\qty{10.0}{g}/\qty{108.0}{g/mol} = \qty{0.0926}{mol}$.

\textbf{5.} Hidrazon $\ce{C9H12N2}$, \qty{148.0}{g/mol}: $\qty{0.0926}{mol} \times \qty{148.0}{g/mol} = \qty{13.7}{g}$.

\textbf{6.} $\ce{C6H5-NH-NH-C(CH3)=CH2}$.

\textbf{7.} Karbon cincin \textit{orto}, karbon \textit{ipso}, kedua nitrogen, karbon hidrazon,
dan $\ce{CH2}$ ujung.

\textbf{8.} Ikatan N--N putus; ikatan C--C terbentuk antara karbon \textit{orto} dan $\ce{CH2}$.

\textbf{9.} Enam elektron, semua komponen \omterm{def:b3:pericyclic:faciality}{suprafasial}: diizinkan secara termal.

\textbf{10.} Sebagai asam Lewis, seng klorida mengkatalisis tautomerisasi, mengikat sebuah nitrogen dan
melemahkan ikatan N--N, serta membantu amonia lepas.

\textbf{11.} Karbon \textit{orto} telah menjadi $sp^3$, dengan sebuah hidrogen: lepasnya proton itu
memulihkan cincin benzena, sebuah gaya dorong yang kuat.

\textbf{12.} $\ce{NH2}$ anilina yang terbentuk menyerang karbon imina, sehingga menutup
cincin beranggota lima (sebuah 2-aminoindolina).

\textbf{13.} Amonia; pembentukan cincin pirola aromatik pada indola mendorongnya.

\textbf{14.} $\ce{C9H12N2 -> C9H9N + NH3}$.

\textbf{15.} \qty{131.0}{g/mol}.

\textbf{16.} $\qty{0.0926}{mol} \times \qty{131.0}{g/mol} = \qty{12.1}{g}$.

\textbf{17.} Pada C3, posisi indola yang paling nukleofilik, yang di sini bebas; serangan di situ
mempertahankan cincin benzena tetap utuh dalam zat antara.

\textbf{18.} $\ce{C6H5NHNH-C(=CH2)-CH2CH3}$ (ujung) dan $\ce{C6H5NHNH-C(CH3)=CH-CH3}$ (dalam).

\textbf{19.} 2-Etilindola dan 2,3-dimetilindola.

\textbf{20.} Yang dalam lebih tersubstitusi; asam kuat menaikkan bagian
yang ujung.

\textbf{21.} 2,3-Dimetilindola, melalui ena-hidrazina yang lebih stabil dan lebih
tersubstitusi.

\textbf{22.} Asetaldehida berkondensasi dengan dirinya sendiri dan rute ena-hidrazina memberikan rendemen
yang buruk; indola dibuat dari hidrazon asam piruvat sebagai gantinya, yang memberikan
asam indola-2-karboksilat, yang kemudian didekarboksilasi.

\textbf{23.} Campuran regioisomer harus dipisahkan, yang mengorbankan rendemen; keton
simetris atau rute yang menetapkan regiokimia lebih disukai.

\textbf{24.} 10.0 g fenilhidrazina dapat memberikan paling banyak \qty{12.1}{g} 2-metilindola.
\end{solution}
